\documentclass[12pt]{article}
\usepackage[utf8]{inputenc}
\usepackage[T1]{fontenc}
\usepackage{hyperref}
\title{Vwanji Language Resources for K-12 Curriculum Development}
\author{Compiled by Claude Sonnet 5 (claude-sonnet-5) for Andr\'as Kornai}
\date{}
\begin{document}
\maketitle

\section{Introduction}
Vwanji (ISO 639-3: wbi; Glottocode vwan1235), also called Wanji, Kivwanji or Kiwanji, is a Bantu language of southwestern Tanzania (Makete District, Njombe Region, and Mbeya Region). Omniglot gives about 28,000 speakers; Joshua Project and ASJP give higher figures. It is written in the Latin script.

\section{Primary Educational Resources}
No primers, textbooks or school programmes for Vwanji turned up in the searches.

\section{Linguistic and Reference Materials}
\begin{itemize}
\item Omniglot page on Vwanji (Ekivwanji): \url{https://omniglot.com/writing/vwanji.htm}
\item Wikipedia entry on the Vwanji language: \url{https://en.wikipedia.org/wiki/Vwanji_language}
\item Journal article, ``A Discourse Analysis of Three Past TAM Forms in Vwanji'': \url{https://journal.fi/store/article/view/68999}
\item ASJP word-list entry for Wanji: \url{https://asjp.clld.org/languages/WANJI}
\end{itemize}

\section{Religious and Community Resources}
\begin{itemize}
\item Scripture Earth page for Vwanji (a Gospel of Luke, 2019 edition, and New Testament audio are reported): \url{https://www.scriptureearth.org/00eng.php?iso=wbi}
\item Joshua Project profile for the Wanji people: \url{https://joshuaproject.net/people_groups/15856}
\end{itemize}

\section{Digital and Audio Resources}
Scripture Earth (above) is the only audio source found.

\section{Curriculum Development Context}
The searches found no school materials or dictionary. The only documentation found is one linguistic article, a short Omniglot overview and scripture portions. A curriculum would need new orthography and teaching materials.
\end{document}